% Source fingerprint: 200c12ed36092f37a69937d328d177bfd44d60fb6439f9b4991dd2197623850a
% T08: raw TeX cells preserved; no numeric highlighting.
\begin{table}[htbp]
\centering\small\renewcommand{\arraystretch}{1.18}\setlength{\tabcolsep}{4pt}
\caption[点值与两种代表]{点值与两种代表。此表区分定义所用的数据，不宣称所有可测函数在每一点都存在这些极限。第28章另证明跳跃处的中点平均公式。}
\label{b1:tab:visual-t08}
\begin{tabularx}{\linewidth}{@{}>{\raggedright\arraybackslash}p{3.1cm}>{\raggedright\arraybackslash}p{4.7cm}>{\raggedright\arraybackslash}X@{}}
\toprule
对象 & 由什么确定 & 点值与相等范围 \\
\midrule
原可测代表 & 给定函数 \(f\) 的逐点取值 & 只改一个零测点也可能破坏该点的近似连续性 \\
有限近似极限 & 误差集合在小球中的相对测度趋零 & 零测修改不改变此极限，但它不必等于原函数点值 \\
近似极限代表 & 用有限近似极限确定 \(\widetilde f\)，其余点按约定补值 & 只有验证与原函数几乎处处相等后，才作为同一等价类的代表 \\
球平均精确代表 & 用球平均极限确定 \(f^*\)，不存在或非有限处补零 & 普通 Lebesgue 点上与近似极限一致；不能在全部异常点逐点认同 \\
Lebesgue 点 & 球平均的 \(|f-f(x)|\) 趋零 & 涉及指定中心值；跳跃点的中点代表不使该点成为 Lebesgue 点 \\
\bottomrule
\end{tabularx}
\end{table}
